\chapter{Introduction}
\label{chap:intro}

Language servers are an essential component of modern development environments, providing rich code intelligence features such as code completion, navigation, error checking, and refactoring support. 
The Language Server Protocol (LSP) \cite{lsp}, introduced by Microsoft, has standardized communication between language servers and editors, allowing a common interface for diverse programming environments. 
Implementing a language server from scratch entails a great deal of boilerplate code and similar, repetitive algorithms for every feature, which makes the work slow; more code also tends to mean more bugs that are harder to identify later.
This thesis aims to make the production of a minimal language server a simpler task by moving as much processing as possible into an agnostic framework built on safe and predictable machinery.
It presents a Haskell framework that, given a Free Foil abstract syntax and a set of functions depending on the language-specific signature, produces a language server communicating via LSP that can store the current environment's configuration, update it on file changes, rename symbols, navigate to a symbol's definition, perform type checking, and present informational, diagnostic, and warning messages.

\section{Binding and Scopes}
\label{sec:bindings-and-scopes}

One of the most persistent aspects of language server implementation, as in compiler implementation, is the correct management of variable bindings and scopes. 
In the following Kotlin example, \kotlin{x} cannot simply be replaced with \kotlin{y} in \kotlin{a}'s body when calling \kotlin{a(y)}, because that would mix the local variable \kotlin{x} with the free \kotlin{y}.

\begin{minted}[fontfamily=txtt]{kotlin}
val y = 1
fun a (x: Int): (Int) -> Int = { y -> x + y }
a(y)
\end{minted}

Handling bound identifiers safely and efficiently is a long-studied problem, with approaches ranging from De Bruijn indices \cite{deBruijn1972} through locally nameless and nominal techniques to intrinsically scoped representations; Chapter~\ref{chap:lr} surveys this design space in detail. 
This thesis builds on Free Foil \cite{freefoil2024}, which combines free scoped monads \cite{freemonads2024} with the type-safe, single-pass approach of Foil \cite{foil2023}. 
Free Foil was chosen because, besides being safe and efficient, it generalizes the Abstract Syntax Tree (AST) over a declarative bifunctor signature, so that features such as navigation to a symbol's definition, renaming, and the reporting of scope-related errors can be implemented once, agnostically, and reused for any language with minimal effort.

Free Foil parameterizes the abstract syntax over a \texttt{binder} type and a signature \texttt{sig} that distinguishes regular nodes from those introducing new scopes, indexing every entity by a phantom scope type \texttt{S} so that the Haskell type system statically rejects operations mixing incompatible scopes. 
Capture-avoiding substitution and normalization thus become available out of the box, with no manual renaming or index shifting during tree traversal. 
Crucially, because the signature is a bifunctor it can implement \texttt{Bifoldable}, enabling a universal recursive traversal of any AST regardless of the specific language. 
By requiring additional type classes only when specific signature data is needed, IDE features can be implemented as generic algorithms, turning language server development from a repetitive, language-specific task into a configuration-driven process: a developer provides only a declarative signature and a minimal set of type class instances, while the framework handles scope management, tree traversal, and LSP communication.

The ambition of describing binding once and deriving editor services generically is not new. 
The most developed prior effort is the line of work on scope graphs and the Spoofax/Statix language workbench \cite{scopeGraphs2015,scopesAsTypes2018,spoofax2010}, which specifies name resolution and type checking declaratively and answers them through a runtime search over an extrinsic graph. 
This thesis instead carries the intrinsic, compiler-checked route of Free Foil across the LSP boundary, trading some of the binding breadth of scope graphs for zero-cost, statically verified scope safety. 
Chapter~\ref{chap:lr} develops this comparison (Section~\ref{subsec:scope-graphs-vs-free-foil}), and Chapter~\ref{chap:conclusion} returns to it.

\section{Contribution}
\label{sec:contribution}

This thesis presents a novel, architecture-driven framework for automating the development of language servers by leveraging the Free Foil library. 
The primary contributions are structured around three core aspects: architectural abstraction, developer ergonomics, and computational guarantees.

First, it introduces a modular, language-agnostic architecture that decouples LSP feature implementation from language-specific syntax. 
By building directly on Free Foil's type-safe scope management and capture-avoiding substitution, the framework abstracts away the repetitive, error-prone logic traditionally required for tree traversal, variable tracking, and scope resolution. 
This enables a single codebase to serve as an agnostic language server capable of handling any language defined via a declarative bifunctor signature.

Second, it defines a lightweight, type-class-driven interface that activates core IDE functionalities with minimal developer overhead. 
Through a carefully designed set of type classes---\texttt{RangedSig}, \texttt{FoldablePat}, \texttt{TypeDeductiveSig}, \texttt{TypedSig}, \texttt{TypedPat}, \texttt{HoverableSig}, \texttt{HoverablePat}, \texttt{TokenizableSig}, \texttt{TokenizablePat}, \texttt{DiagnosableSig}, \texttt{DiagnosablePat}---the framework exposes syntax highlighting, diagnostics, go-to-definition, symbol renaming, hover tooltips, and basic type checking. 
A language integrator implements only a declarative AST signature and a small set of type class instances: 14 type-class methods together with the \texttt{buildAsts} function, 15 functions in total. 
All scope management, recursive traversal, and LSP communication are handled automatically by the library, shifting the development paradigm from boilerplate-heavy implementation to configuration-driven setup.

Third, it provides a formal worst-case complexity analysis of all server algorithms, demonstrating that every operation runs in at most linear time relative to the AST size ($O(K)$); position-based queries such as go-to-definition and hover descend a single path and, for fixed-arity grammars, run in time linear in the tree depth ($O(L)$, with $L \leq K$). This guarantees predictable performance and scalability for typical development workflows. 

The framework is validated by two independent example language servers built on it without modifying any core module: LambdaPi, a dependently typed lambda calculus with $\Pi$-types, and Flan, a functional language with let-bindings, pair patterns, and conditionals that exercises the full feature set. 
From a practical standpoint, the proposed framework significantly accelerates the creation of language servers, making it particularly valuable for educational projects, rapid language prototyping, and the integration of domain-specific languages into modern IDEs. 
The thesis also critically evaluates the current implementation, identifying key limitations such as the absence of global namespace resolution, reliance on simple annotation-based type inference (lacking constraint unification), and bulk cache updates upon file modification. 
In response, it outlines concrete directions for future work, including incremental type checking, interval-map indexing for logarithmic-time ($O(\log K)$) position lookups, single-pass tree traversals, and generic default implementations for common LSP features. 
The source code and documentation are provided to serve as a foundational toolkit for further research in automated development tooling.

\section{Thesis Outline}
\label{sec:outline}

The remainder of this thesis is organized as follows. 
Chapter~\ref{chap:lr} reviews the techniques for safe and efficient handling of bound variables, the existing language-agnostic models of name binding and resolution, and the Language Server Protocol. 
Chapter~\ref{chap:met} presents the methodology: the design requirements, the agnostic tree-traversal strategy, and the algorithms behind each supported feature. 
Chapter~\ref{chap:impl} documents the implementation, covering the LSP interface, caching, name-collision handling, and the Flan example integration. 
Chapter~\ref{chap:eval} evaluates the framework against its design criteria and analyzes its automation and computational complexity, together with current limitations and directions for future work. 
Chapter~\ref{chap:conclusion} concludes.
